% -----------------------------------------------------------------------------
% This file is automatically generated.
% Do not edit manually.
% Source: notebooks/support/hypothesis-testing/support.Rmd
% -----------------------------------------------------------------------------

\noindent The sample size in the \emph{Case: European innovation subsidies} depends on the critical region chosen. When the null hypothesis $H_0 : M \geq 120,000$ is rejected when the sample yields no errors, the critical region is $\{k | k = 0\}$. We first create an \ttblue{mus\_obj} object, and load it with the parameters.
\par\addvspace{\topsep}
\begingroup
\fvset{listparameters={%
  \setlength{\topsep}{0pt}%
  \setlength{\partopsep}{0pt}%
  \setlength{\parsep}{0pt}%
  \setlength{\itemsep}{0pt}%
}}
\begin{Verbatim}[breaklines=true,formatcom=\color{ada_blue}]
subsidies <- mus_obj(
  cl = 0.95,
  popBv = 12000000,
  pm = 120000
)
subsidies <- size(subsidies,
  ee = 0
)
subsidies$n
\end{Verbatim}
\begin{Verbatim}[breaklines=true,formatcom=\color{ada_light_blue}]
[1] 299
\end{Verbatim}
\endgroup
\par\addvspace{\topsep}
\noindent This result is exactly equal to that of the fixed-attribute sample:
